\scanpage{79}
dont ils forment une subdivision cellulaire -- ce cercle définit une structure réelle canonique sur la droite projective complexe $\Sigma$, noté $\Sigma(\mathbb R)$, \deleted{\unsure{ou} $\Sigma_{\mathbb R}$}\inserted{$\Sigma_{\mathbb R}$}. De plus, il y a une métrique riemannienne homogène unique \inserted{(à facteur const-près)} sur $\Sigma$ invariante par $D$ (c'est pour elle que $\Sigma(\mathbb R)$ est un \uline{grand} cercle).
\margin{il est donc \deleted{\unsure{défini}}}

$\Sigma_{\mathbb R}\deleted{(\mathbb R)}$ partage $\Sigma$ en deux hémisphères, \deleted{et} en correspondance biunivoque avec $\underline\omega$ par la correspondance de Stokes
\[
(49)\qquad \pi_0(\Sigma\setminus\Sigma_{\mathbb R})\simeq\underline\omega.
\]
On a
\[
\Sigma^D=\Sigma^{D^+}=\Sigma^{\underline\omega}=\Sigma^{\{\omega\}}\unsure{\subsetneq|\subset}\Sigma\setminus\Sigma_{\mathbb R}
\qquad(\omega\in\underline\omega)
\]
et $\operatorname{card}\Sigma^D=2$, de façon précise, chacun des hémisphères ouverts contient un pt et un seul de $\Sigma^D$, d'où
\[
(50)\qquad \Sigma^D\simeq\underline\omega,
\]
on désigne par $P_\omega$ le point de $\Sigma^D$ qui correspond à $\omega\in\underline\omega$. Il est caractérisé par le fait que l'opération de $\omega$ dans l'espace tangent à $\Sigma$ en $P_\omega$ est une rotation d'angle $+2\pi/n$ (alors qu'en $P_{\omega'}$ elle est d'angle $-2\pi/n$). Pour $s\in S$, on désigne par $R_s$ le point correspondant de $\Sigma_{\mathbb R}$, pour $a\in A$ on désigne par $Q_a$ le centre de l'arête $a$.
\margin{78}
